% !TeX root = ../../Book2.tex
% 纲要标题：### 3.3 正合列
% 纲要位置：计划纲要.md，第 899 行。

\section{正合列}
\label{b2:sec:0303}

一个模同态的核说明哪些元素被送到零，像说明哪些元素能够到达。若把两个同态首尾相接，便可以进一步问：后一个同态送到零的元素，是否恰好都是前一个同态的像？正合性把这种关系写成一个等式。本节及下一节中的模都是同一个含幺环 \(R\) 上的幺性左模，箭头都是 \(R\) 模同态；不要求 \(R\) 交换。

% 纲要内容（仅作写作计划，尚非教材正文）：
%
% * 短正合列与核—余核描述
% * 分裂短正合列
% * 正合性在交换图中的使用
%

\subsection{短正合列与核—余核描述}
\label{b2:subsec:0303:01}

\begin{definition}\label{b2:def:exact-sequence}
对于两个相接的同态 \(A\xrightarrow{u}B\xrightarrow{v}C\)，若
\[
\operatorname{im}u=\ker v,
\]
就称这段序列在 \(B\) 处\term[zhenghe]{正合}[exact]。一列相接的模同态若在每个同时有入箭头与出箭头的位置都正合，就称为\term[zhenghe-lie]{正合列}[exact sequence]。特别地，正合列
\begin{equation}\label{b2:eq:short-exact-sequence}
0\longrightarrow A\xrightarrow{i}B\xrightarrow{p}C\longrightarrow0
\end{equation}
称为\term[duan-zhenghe-lie]{短正合列}[short exact sequence]。
\end{definition}

这里的 \(0\) 表示零模；进出零模的同态都只能是零映射。因此，\eqref{b2:eq:short-exact-sequence}在 \(A\) 处正合表示 \(i\) 单射，在 \(C\) 处正合表示 \(p\) 满射，在 \(B\) 处正合则表示 \(\operatorname{im}i=\ker p\)。短正合列同时记住了子模怎样嵌入以及商映射怎样给出，并不只记三个模的同构类型。

\ProseParagraphBreak
正合性蕴含 \(vu=0\)，但只验证这个复合为零还不够。例如 \(\mathbb Z\xrightarrow{\times2}\mathbb Z\xrightarrow{0}\mathbb Z\) 中，复合为零，而中间位置的像为 \(2\mathbb Z\)、核为 \(\mathbb Z\)，二者不等。正合性还要求每个被 \(v\) 送到零的元素都能写成 \(u(a)\)，这个“能够找到原像”的条件会成为图追踪的依据。

\ProseParagraphBreak
对每个整数 \(n\ge2\)，有短正合列
\begin{equation}\label{b2:eq:integer-short-exact}
0\longrightarrow\mathbb Z\xrightarrow{\times n}\mathbb Z
\xrightarrow{q}\mathbb Z/n\mathbb Z\longrightarrow0,
\qquad q(a)=a+n\mathbb Z.
\end{equation}
乘以 \(n\) 在整数中单射，\(q\) 满射，而 \(\ker q=n\mathbb Z\) 恰为前一箭头的像。对任意子模 \(A\le B\)，其包含映射及商映射同样给出 \(0\to A\to B\to B/A\to0\)。这两种例子的标量环都是 \(\mathbb Z\) 时，正合性就是交换群同态的相应性质。

\begin{proposition}[短正合列中的核与余核]\label{b2:prop:short-exact-kernel-cokernel}
设 \eqref{b2:eq:short-exact-sequence}短正合。则 \(i\) 给出 \(A\cong\ker p\)，\(p\) 给出同构
\[
B/\operatorname{im}i\longrightarrow C,
\qquad b+\operatorname{im}i\longmapsto p(b).
\]
更具体地，对任意模 \(X\)，有以下两项性质：
\begin{enumerate}
\item 若同态 \(f:X\to B\) 满足 \(pf=0\)，则存在唯一同态 \(\widetilde f:X\to A\)，使 \(i\widetilde f=f\)。
\item 若同态 \(g:B\to X\) 满足 \(gi=0\)，则存在唯一同态 \(\widetilde g:C\to X\)，使 \(\widetilde g p=g\)。
\end{enumerate}
\end{proposition}

\begin{proof}
\(i\) 单射且像为 \(\ker p\)，所以限制值域后就是第一个同构。\(p\) 满射且核为 \(\operatorname{im}i\)，由\cref{b2:thm:module-first-isomorphism}得到第二个同构。

若 \(pf=0\)，则对每个 \(x\in X\)，\(f(x)\in\ker p=\operatorname{im}i\)。因为 \(i\) 单射，存在唯一的 \(a\in A\) 满足 \(i(a)=f(x)\)，令 \(\widetilde f(x)=a\)。在这个等式中代入和或左标量倍，再用 \(i\) 的单射性，便得 \(\widetilde f\) 线性，其取值也已被等式唯一确定。

若 \(gi=0\)，则 \(g\) 在 \(\operatorname{im}i\) 上为零。由\cref{b2:prop:quotient-module}，它唯一地经 \(B/\operatorname{im}i\) 分解，再通过上面的同构视为 \(C\) 上的同态。等价地，定义 \(\widetilde g(p(b))=g(b)\)；两个原像之差在 \(\ker p=\operatorname{im}i\) 中，保证取值相同，而 \(p\) 满射保证唯一性。
\end{proof}

前一项性质把 \(i\) 描述为 \(p\) 的核映射，后一项把 \(p\) 描述为 \(i\) 的余核映射。这里所说的核与余核都带着指定的同态，而不是把 \(A\) 与某个核、\(C\) 与某个商模任意认同。

\ProseParagraphBreak
一般同态 \(f:M\to N\) 可以放进正合列
\[
0\longrightarrow\ker f\longrightarrow M\xrightarrow{f}N
\longrightarrow\operatorname{coker}f\longrightarrow0.
\]
若在像处把它分开，就得到两个短正合列
\[
0\longrightarrow\ker f\longrightarrow M\longrightarrow\operatorname{im}f\longrightarrow0,
\]
\[
0\longrightarrow\operatorname{im}f\longrightarrow N\longrightarrow\operatorname{coker}f\longrightarrow0.
\]
第一列把 \(f\) 的值域限制为其像，第二列使用像在 \(N\) 中的包含映射。每处正合性分别来自核、像及商模的定义。这说明一个同态的核与余核，怎样通过两个短正合列与原来的模联系起来。

\subsection{分裂短正合列}
\label{b2:subsec:0303:02}

对任意两个模 \(A,C\)，坐标嵌入和投影给出短正合列
\begin{equation}\label{b2:eq:standard-split-sequence}
0\longrightarrow A\xrightarrow{\iota_A}A\oplus C
\xrightarrow{\pi_C}C\longrightarrow0,
\quad \iota_A(a)=(a,0),\quad \pi_C(a,c)=c.
\end{equation}
这个例子有两个额外的同态：\((a,c)\mapsto a\) 可以收回 \(A\) 分量，\(c\mapsto(0,c)\) 可以为每个 \(c\) 选择一个与加法及标量作用相容的原像。一般短正合列不一定有这样的选择。

\begin{definition}\label{b2:def:section-retraction}
对同态 \(p:B\to C\)，满足 \(ps=\operatorname{id}_C\) 的同态 \(s:C\to B\) 称为 \(p\) 的一个\term[jie-mian]{截面}[section]。对同态 \(i:A\to B\)，满足 \(ri=\operatorname{id}_A\) 的同态 \(r:B\to A\) 称为 \(i\) 的一个\term[suohui]{缩回}[retraction]。这两种等式分别说 \(s\) 是 \(p\) 的右逆、\(r\) 是 \(i\) 的左逆。
\end{definition}

\begin{theorem}[短正合列的分裂]\label{b2:thm:split-exact-sequence}
对于短正合列 \eqref{b2:eq:short-exact-sequence}，以下条件等价：
\begin{equivalences}
\item\label{b2:item:split-section} \(p\) 有截面 \(s:C\to B\)。
\item\label{b2:item:split-retraction} \(i\) 有缩回 \(r:B\to A\)。
\item\label{b2:item:split-compatible-sum} 存在同构 \(\theta:A\oplus C\to B\)，使
\(\theta\iota_A=i\) 且 \(p\theta=\pi_C\)。
\end{equivalences}
满足这些条件的短正合列称为\term[fenlie-duan-zhenghe-lie]{分裂短正合列}[split short exact sequence]。
\end{theorem}

\begin{proof}
先由\itemref{b2:item:split-section}证明\itemref{b2:item:split-compatible-sum}。给定截面 \(s\)，定义
\[
\theta(a,c)=i(a)+s(c).
\]
它是模同态，且满足两项相容等式。若 \(\theta(a,c)=0\)，施加 \(p\) 得 \(c=0\)，再由 \(i\) 单射得 \(a=0\)。对每个 \(b\in B\)，有
\[
p\bigl(b-s(p(b))\bigr)=0.
\]
由正合性存在 \(a\in A\)，使 \(b-s(p(b))=i(a)\)，所以 \(b=\theta(a,p(b))\)。因此 \(\theta\) 单射且满射，是所求同构。

由\itemref{b2:item:split-compatible-sum}得到\itemref{b2:item:split-retraction}，只需令 \(r=\pi_A\theta^{-1}\)。于是
\(ri=\pi_A\theta^{-1}\theta\iota_A=\operatorname{id}_A\)。

最后由\itemref{b2:item:split-retraction}构造\itemref{b2:item:split-section}。对 \(c\in C\)，用 \(p\) 的满射性选 \(b\in B\) 满足 \(p(b)=c\)，令
\[
s(c)=b-i(r(b)).
\]
若换成另一原像 \(b'\)，则 \(b'-b=i(a)\) 对某个 \(a\in A\) 成立，故
\[
\bigl(b'-i(r(b'))\bigr)-\bigl(b-i(r(b))\bigr)
=i(a)-i(ri(a))=0.
\]
因此 \(s\) 与原像选择无关。检验和与左标量倍时，可以分别选取 \(b_1+b_2\) 及 \(xb\) 为原像，其中 \(x\in R\)；定义立即给出 \(s(c_1+c_2)=s(c_1)+s(c_2)\) 及 \(s(xc)=xs(c)\)。又有 \(ps(c)=p(b)=c\)，所以 \(s\) 是截面。
\end{proof}

通过上述构造相配的 \(r,s\) 还满足
\[
ri=\operatorname{id}_A,\qquad ps=\operatorname{id}_C,
\qquad rs=0,\qquad ir+sp=\operatorname{id}_B.
\]
这些等式分别说明两个分量怎样取出并重新合成。特别地，\(B=i(A)\oplus s(C)\)，这里是\cref{b2:prop:internal-direct-sum}中的内直和。定义中的相容条件不可省略：判断给定短正合列是否分裂，需要同构尊重已经指定的 \(i,p\)，不能只写一个没有说明映射的 \(B\cong A\oplus C\)。

\begin{corollary}[商模自由时的分裂]\label{b2:cor:free-quotient-splits}
若短正合列 \(0\to A\to B\xrightarrow{p}F\to0\) 的商模 \(F\) 自由，则该列分裂。
\end{corollary}

\begin{proof}
在\cref{b2:prop:free-lifting}中，取满同态为 \(p:B\to F\)，取从自由模出发的同态为 \(\operatorname{id}_F:F\to F\)。该命题给出同态 \(s:F\to B\)，使 \(ps=\operatorname{id}_F\)。这就是 \(p\) 的截面，再用\cref{b2:thm:split-exact-sequence}即可。
\end{proof}

在域上，每个向量空间都有基，因此向量空间的每个短正合列都分裂。对一般环，这个结论会失败。\eqref{b2:eq:integer-short-exact}若有截面 \(s:\mathbb Z/n\mathbb Z\to\mathbb Z\)，则
\[
n\,s(1+n\mathbb Z)=s\bigl(n(1+n\mathbb Z)\bigr)=0.
\]
整数中没有非零的有限阶元素，故 \(s(1+n\mathbb Z)=0\)，不可能满足 \(qs(1+n\mathbb Z)=1+n\mathbb Z\)。因此这列不分裂。模论中“为每个商类选择一个代表元”与“选择一个模同态截面”是两件不同的事，后者必须保留代表元之间的加法与标量关系。

\ProseParagraphBreak
即使分裂存在，截面也通常不唯一。若 \(s\) 是一个截面，\(h:C\to A\) 是任意同态，则 \(s+ih\) 仍是截面，因为 \(pi=0\)。反过来，两个截面之差被 \(p\) 零化，由\cref{b2:prop:short-exact-kernel-cokernel}唯一地写成 \(ih\)。所以分裂给出一种直和分解；改变截面，可能改变作为第二个分量嵌入 \(B\) 的子模。

\subsection{正合性在交换图中的使用}
\label{b2:subsec:0303:03}

若图中任意两条起点、终点相同的有向路径，其同态复合都相同，就称这个图为\term[jiaohuan-tu]{交换图}[commutative diagram]。下图的两行都是短正合列：
\begin{equation}\label{b2:eq:short-exact-diagram}
\begin{tikzcd}[column sep=1.6em,row sep=2em]
0\arrow[r]&A\arrow[r,"i"]\arrow[d,"\alpha"']&B\arrow[r,"p"]\arrow[d,"\beta"]&C\arrow[r]\arrow[d,"\gamma"]&0\\
0\arrow[r]&A'\arrow[r,"i'"']&B'\arrow[r,"p'"']&C'\arrow[r]&0.
\end{tikzcd}
\end{equation}
交换条件正是 \(\beta i=i'\alpha\) 和 \(p'\beta=\gamma p\)。它使沿上行再向下与先向下再沿下行得到同一个元素；两行的正合性则允许我们在某些元素的像为零时，沿前一个箭头寻找原像。

\begin{lemma}[短五引理]\label{b2:lem:short-five}
在\eqref{b2:eq:short-exact-diagram}中，若 \(\alpha,\gamma\) 都单射，则 \(\beta\) 单射；若 \(\alpha,\gamma\) 都满射，则 \(\beta\) 满射。特别地，若 \(\alpha,\gamma\) 都是同构，则 \(\beta\) 也是同构。
\end{lemma}

\begin{proof}
先假设 \(\alpha,\gamma\) 单射。若 \(\beta(b)=0\)，交换性给出 \(\gamma(p(b))=p'(\beta(b))=0\)，由 \(\gamma\) 单射得 \(p(b)=0\)。上行正合，故可取 \(a\in A\)，使 \(b=i(a)\)。于是 \(i'(\alpha(a))=\beta(i(a))=0\)。\(i'\) 及 \(\alpha\) 都单射，所以 \(a=0\)，进而 \(b=0\)。

再假设 \(\alpha,\gamma\) 满射。给定 \(b'\in B'\)，先利用 \(\gamma\) 满射选 \(c\in C\)，使 \(\gamma(c)=p'(b')\)；再利用 \(p\) 满射选 \(b\in B\)，使 \(p(b)=c\)。这时未必已有 \(\beta(b)=b'\)，但
\[
p'\bigl(b'-\beta(b)\bigr)=p'(b')-\gamma(p(b))=0.
\]
下行正合，故差值可写成 \(i'(a')\)。由 \(\alpha\) 满射选 \(a\in A\)，使 \(\alpha(a)=a'\)，便有
\[
\beta\bigl(b+i(a)\bigr)=\beta(b)+i'(\alpha(a))=b'.
\]
因此每个 \(b'\) 都有原像，\(\beta\) 满射。
\end{proof}

这一结果称为\term[duan-wu-yinli]{短五引理}[short five lemma]。证明中的第一次提升只保证商中的像正确；正合性把剩余差值放回左端，再用左端的映射修正。这种“先满足一个条件，再修正差值”的做法，比同时猜出三个模之间的全部对应更容易操作。

\ProseParagraphBreak
引理要求先有交换图。两个短正合列即使左右两端分别相同，也不一定能找到中间的相容同构。例如
\[
0\longrightarrow\mathbb Z/2\mathbb Z\longrightarrow\mathbb Z/4\mathbb Z
\longrightarrow\mathbb Z/2\mathbb Z\longrightarrow0
\]
的左箭头把 \(1+2\mathbb Z\) 送到 \(2+4\mathbb Z\)，而将中间项换成 \((\mathbb Z/2\mathbb Z)^2\)、两箭头换成第一坐标嵌入与第二坐标投影，又得到一个短正合列。前者中间项含四阶元素，后者每个元素都被二零化，故两者不可能同构。短五引理不能仅从端点相同推出中间项相同；它说明的是一组已经相容的映射何时必为同构。
